\documentclass[12pt, letterpaper]{article}

% Essential packages
\usepackage[utf8]{inputenc}
\usepackage[T1]{fontenc}
\usepackage{geometry}
\usepackage{setspace}
\usepackage{parskip} % Handle paragraph spacing
\usepackage{titlesec} % For customizing section titles

% Page layout settings
\geometry{
    top=1in,
    bottom=1in,
    left=1in,
    right=1in
}

% Line spacing
\setstretch{1.15}

% Remove page numbers for a clean look (optional, comment out if you want numbers)
\pagestyle{empty}

\begin{document}

% --- HEADER START ---
\begin{center}
    \textbf{\Large Kritarth Dandapat } \\ \vspace{4pt}
    Computer Science $\cdot$ University at Buffalo \\
    kritarth@buffalo.edu $\cdot$ +1 716 612 0016 \\ \vspace{4pt}
    % Optional: Link to LinkedIn or Portfolio
    % \text{linkedin.com/in/yourprofile} 
\end{center}

\vspace{0.5em}
\hrule
\vspace{1.5em}

\begin{center}
    \textbf{\Large PERSONAL STATEMENT}
\end{center}
\vspace{1em}
% --- HEADER END ---

At nineteen years of age, I am applying for a PhD program, since I have always maintained an \textbf{accelerated pace} throughout my academic journey. I finished high school a year ahead, and I am on track to complete my college again a year ahead of schedule. The desire to achieve such momentum is not simply a matter of speed or rush, but from an intense intellectual hunger to reach the \textbf{frontier of research} as soon as possible. This trajectory has taught me \textbf{resilience and problem-solving traits} that were first tested when I transitioned from the Indian educational system to that of the United States.

Growing up in India, my high school years were calibrated entirely toward one goal: to secure a higher rank on the IIT-JEE (Advanced) examination. When I enrolled for FIITJEE, a premier coaching Institute for IIT-JEE, along with my school curriculum during my Class seven through twelve, I had no space for additional activities. When I pivoted to the U.S. university system, I had virtually no lead time---only one month to prepare for the SATs, IELTS, and applications. That month was a crash course in \textbf{resourcefulness}. I learned to maximize limited time to achieve a specific outcome, a skill that has defined my undergraduate career.

In college, I have consistently been allowed to opt for \textbf{21 or 22 credits per semester}, significantly exceeding the standard load, based on my being an honors student and holding a high GPA. I did not do this to check boxes, but I refused to be limited by a standard timeline. I sought out the most challenging coursework available, enrolling in advanced classes like \textbf{Distributed Systems and Quantum Computing} well before my peers.

However, this aggressive pursuit of knowledge taught me a vital lesson about balance and prioritization. In the Spring of 2025, I faced my most demanding semester, balancing 22 credits of difficult coursework with research and tutoring. For the first time, I received a `B' grade. While my GPA dipped slightly, I view that semester as a triumph of challenges. I could have stepped back from my research or stopped tutoring to secure a perfect transcript, but I chose to honor my commitments to my lab and my students. This choice proved to me that the \textbf{quality of my contribution to science} and my community matters more than a flawless academic record.

My \textbf{commitment to research} has been the central pillar of my time at university. For my entire second year, I worked under Professor Wenyao Xu while simultaneously tutoring. During the summer break, I expanded my scope by taking on a second research project with Professor Jiyau Peng and continuing during fall 2025. While managing two research streams was unconventional, I viewed these years as my opportunity to grasp as much as possible.

This effort has yielded tangible results. When I began with Professor Xu, I had no prior knowledge of the specific domain; I \textbf{taught myself mobile application development and AI/ML} on the job. My rapid progress led to me being paid within two months---a rarity for undergraduates in that lab. Similarly, my work with Professor Peng pushed me to new heights. I \textbf{published a paper} and, with his support, received the \textbf{PEARL Award}, a prestigious \$2,500 grant awarded only once during an undergraduate career. My CV projects this strong work ethic, demonstrating that I have the discipline required for doctoral study.

These experiences have shaped how I should contribute to an academic community. Being consistently younger than my cohort has required me to be adaptable, learning to listen and \textbf{collaborate effectively} with those more experienced than myself. Additionally, my dedication to tutoring reflects my belief in mutual support.

I understand that there will be multiple challenges while navigating through complex educational systems, and as a prospective PhD student, I hope I will be able to bring such \textbf{resilience, collaborative spirit, and dedication} to my future research group.

\end{document}